\documentclass{article} % say
\usepackage{amsmath}
\usepackage{amssymb}
\usepackage{gensymb}
\usepackage{hyperref}
\usepackage{tikz}
\usepackage{tkz-euclide}
\usepackage{listings}
\usetikzlibrary {babel,angles,quotes,positioning}
\usepackage[utf8]{inputenc}
\usepackage[russian]{babel}

\usepackage{mathtools}
\DeclarePairedDelimiter\ceil{\lceil}{\rceil}
\DeclarePairedDelimiter\floor{\lfloor}{\rfloor}
\begin{document}

\textbf{Задача}: Дано множество $S$ целых чисел и целое неотрицательное целое число $c$. Нужно найти сколько существует пар неотрицательных целых чисел $x$ и $y$, таких что $0 \leqslant x \leqslant y \leqslant c$ и что $x+y$ и $y-x$ не принадлежат $S$. $S$ задано в отсортированном порядке. \textbf{Задача отличается от оригинальной тем, что числа из $S$ не ограничены $c$}.

\clearpage

\textbf{Решение}: 

Посчитаем количество таких пар $x$ и $y$, что они лежат в нужных границах и:
\begin{enumerate}
  \item $x+y$ принадлежит $S$ ($N^{+}$),
  \item $y-x$ принадлежит $S$ ($N^{-}$),
  \item оба $y \pm x$ принадлежат $S$ ($N^{\pm}$).
\end{enumerate}

Тогда ответ это: $\frac{(c + 1) (c + 2)}{2} - N^{+} - N^{-} + N^{\pm}$ согласно принципу включения-исключения.

Для подсчёта $N^{+}$ сделаем следующее. Переберём число $a$ из $S$. Какие $x$ могут давать $a$ в сумме? Они должны быть больше $0$ и $a - c$ (иначе не хватит $y$, чтоб добить до суммы). С другой стороны они должны быть меньше $c$ и $\lfloor \frac{a}{2} \rfloor$ (чтоб $y$ был не меньше $x$). Все числа из полученного интервала задают $x$, а значит и $y$. Эту часть можно сделать за $\mathcal{O}(|S|)$.

Аналогично подсчитаем $N^{-}$. Тут для числа $a$ будем перебирать $y$. Он может быть любым от $a$ до $c$. Сложность такая же.

Наконец, рассмотрим $N^{\pm}$. Заметим, что $y \pm x$ одной чётности, поэтому задача разбивается на две одинаковые: одна для чётных чисел из $S$, другая для нечётных. Обозначим $a = y - x$, $b = y + x$. Зафиксируем $a$, найдём все $b$, такие что существуют $x, y$, которые удовлетворяют условию задачи. Во-первых, $b$ должно быть не меньше $a$. Во-вторых, оно не может быть больше, чем $2c - a$ (иначе $x$ или $y$ должен быть больше $c$). Все числа, которые удовлетворяют этим условиям, подходят. Так как $2c - a$ убывает при росте $a$, то можно найти все подходящие $b$ для всех $a$ в порядке возрастания с помощью метода скользящего окна за $\mathcal{O}(|S|)$.  

\lstinputlisting[language=C++]{main.cpp}

\textbf{Источник}: \url{https://codeforces.com/contest/1935/problem/D}

\end{document}


